\section{Introduction}
\label{sec:introduction}

\subsection{The Calibration Problem in Redistricting Ensembles}

The algorithmic redistricting toolkit has matured to the point where
practitioners can generate thousands of valid redistricting plans in minutes.
But generating plans is not the same as sampling from the correct
distribution.
For statistical inference --- answering questions like ``what fraction of all
valid North Carolina plans produce at least 7 Democratic congressional seats?''
--- the distribution of the generated ensemble matters as much as its size.

Two failure modes are common in practice.
First, optimisation methods such as \cs{}~\citep{b11paper} and
\textsc{ShortBurst}~\citep{cannon2022,g6paper} do not target any
distribution: they search for plans that optimise an objective (edge-cut
minimality, seat outcome), not plans drawn uniformly at random.
An optimisation ensemble answers ``what are the best plans?'', not ``what
fraction of all plans are Democratic-leaning?''

Second, Markov chain methods such as \recom{}~\citep{deford2021} target a
specific stationary distribution, but the chain must \emph{mix} before the
sample approximates that distribution.
For large states with many districts, mixing is not guaranteed in
practically feasible chain lengths.
\recom{}-based $p$-value claims require that the chain has mixed, which is
difficult to certify and is often left unverified in published analyses.

Sequential Monte Carlo (\smc{}) circumvents both failure modes.
\smc{} constructs each plan sequentially, one district at a time, and
importance-weights each plan by the product of the inverse proposal densities
at each stage~\citep{fifield2020}.
The result is a weighted sample that is \emph{asymptotically correct} for the
uniform distribution over valid $k$-district plans, without any Markov chain
mixing assumption.

\subsection{Three Contributions}

This paper makes three contributions.

\paragraph{Contribution 1: The \texttt{bisect-smc} Rust implementation.}
We implement \smc{} as a Rust crate (\texttt{bisect-smc}) with L0, L1, and
ignored real-data L2 tests,
Rayon parallelism at the per-stage particle-proposal step, and a full SHA-256
audit chain.
The implementation is exposed to the CLI as
\texttt{bisect ensemble --method smc --state NC --year 2020 --particles 5000 --base-seed 42}.

\paragraph{Contribution 2: The importance-weighting derivation.}
We derive the weight increment formula $\Delta\log w_i = \log|{\rm
valid\_cuts}(T)|$ from first principles (Section~\ref{sec:algorithm}), making
explicit why the uniform spanning tree proposal requires this correction to
recover the uniform distribution over valid plans.

\paragraph{Contribution 3: The largest-component seeding rule.}
We prove that seeding each stage from the largest connected component of
unassigned tracts --- rather than from an arbitrary tract --- is necessary to
maintain the connectivity invariant that the unassigned subgraph remains
connected after each district assignment
(Section~\ref{subsec:connectivity-invariant}).
Without this rule, the algorithm may paint itself into a corner where
remaining tracts are isolated, making it impossible to form the final $k$-th
district.

\subsection{Why Calibration Matters for Redistricting}

The canonical inference question in redistricting is the percentile claim:
``the enacted plan falls at the Xth percentile of valid plans by partisan
outcome.''
This claim has legal traction --- courts can assess whether the enacted plan
is an extreme outlier or a typical member of the valid-plan space --- but only
if ``the Xth percentile'' refers to the correct distribution.

If the ensemble is biased toward compact plans (because \recom{} oversamples
compact configurations relative to the uniform distribution), a plan at the
90th percentile of that biased ensemble may be at the 60th percentile of the
true uniform distribution.
This is not a theoretical concern: \citet{fifield2020} demonstrate empirically
that ReCom-based ensembles for small synthetic graphs deviate significantly from
the uniform distribution, even after chains that appear to have mixed by
conventional diagnostics.

\smc{} addresses this concern by construction when the proposal support and
weighting model match the declared target distribution.
For any bounded function $f$ of redistricting plans, the \smc{} estimator
$\hat{f}_N = \sum_i w_i f(\pi_i)$ converges to
$\mathrm{E}_\pi[f]$ (the expectation under the uniform distribution) as $N \to
\infty$, without Markov chain assumptions.

\subsection{Scope and Limitations}

This paper describes \smc{} as implemented in \texttt{bisect-smc} and
validated on synthetic graphs, North Carolina and Wisconsin real-data L2 tests,
and Phase~2 large-$k$ Texas/California runs reported in
Section~\ref{sec:phase2}. \texttt{SeedCompositor::SmcPercentile} is now
implemented for selected-plan export from the weighted ensemble. The remaining
claim boundary is external replication: the paper reports local implementation
evidence, not an independent court-ready validation against every outside SMC
implementation.

\subsection{Paper Structure}

Section~\ref{sec:algorithm} gives the formal \smc{} algorithm, including the
seed schedule, importance weights, and systematic resampling step.
Section~\ref{sec:comparison} compares \smc{} with the four other redistricting
methods in the toolkit.
Section~\ref{sec:audit} describes the audit chain and 5-step verification
protocol.
Section~\ref{sec:implementation} covers the Rust implementation details:
Rayon parallelism, Kahan summation, and memory layout.
Section~\ref{sec:conclusion} concludes with deployment guidance and open
questions.
